\section{Adaptive algorithm, cache batches, and cost}
\label{sec:adaptive}
%======================================================================

The local indicator $\eta_H$ is inexpensive and is evaluated at every
coarse step.  The exact certificate requires two operator-norm calculations
and is therefore evaluated only at selected checkpoints.  The fine mesh is
not a reference target: it is an adaptive computational space whose
adequacy is tested by the total corrector defects.

%----------------------------------------------------------------------
\subsection{Refinement logic}
\label{sec:adaptive-logic}
%----------------------------------------------------------------------

A \emph{cache batch} is a maximal sequence of coarse adaptive steps during
which $\mathcal T_h$, $\ell$, and $\kappa$ remain fixed.  Coarse refinement
may still alter $I_H$ and some patch geometries, but local corrector, Riesz,
and mixed factors are retained on all unaffected patches.  A batch ends
when $\mathcal T_h$ is refined or $\ell$ is changed.  This terminology is
purely computational and introduces neither a reference solution nor an
epoch-switch criterion.

To avoid an exact check after every solve, store the last checkpoint index
$n_{\mathrm{cert}}$ and the corresponding value
$\eta_H^{\mathrm{cert}}$.  A new check is triggered when
\begin{equation}
\label{eq:cert-trigger-reduction}
  \eta_H^{(n)}\leq q_{\mathrm{cert}}\eta_H^{\mathrm{cert}},
  \qquad0<q_{\mathrm{cert}}<1,
\end{equation}
or
\begin{equation}
\label{eq:cert-trigger-delay}
  n-n_{\mathrm{cert}}\geq m_{\mathrm{cert}}.
\end{equation}
A check is also forced before termination and after any change of
$\mathcal T_h$ or $\ell$.

When a total-defect check fails, the localization quantities in
\eqref{eq:localization-certificate} separate the two possible causes.
If the certified discrete localization defect exceeds a prescribed
$\tau_{\mathrm{loc}}$, increase $\ell$.  Otherwise refine $\mathcal T_h$
with the local contributions generated by the dominant generalized
eigenvectors in \eqref{eq:defect-eigenproblems}.  If
$v_H^{\mathrm p}$ and $v_H^{\mathrm a}$ are the current dominant modes,
write their patchwise majorants as sums of element contributions and mark
with
\begin{equation}
\label{eq:fine-defect-indicator}
  (\eta_{h,K}^{\mathrm{def}})^2
  :=(\eta_{\mathrm a,K}^{\mathrm{tot}})^2
    +\omega_{\mathrm p}(\eta_{\mathrm p,K}^{\mathrm{tot}})^2,
  \qquad0<\omega_{\mathrm p}\leq1.
\end{equation}
The adjoint part directly enters reliability; the primal part is needed
for the stability transfer.  Taking $\omega_{\mathrm p}=1$ is the safe
default, while a smaller value is reasonable once $\betaInf$ is well
separated from zero.  Choose a fine-mesh D\"orfler set
$\mathcal M_h\subseteq\mathcal T_h$ such that
\begin{equation}
\label{eq:fine-defect-dorfler}
  \sum_{K\in\mathcal M_h}(\eta_{h,K}^{\mathrm{def}})^2
  \geq\theta_h
       \sum_{K\in\mathcal T_h}(\eta_{h,K}^{\mathrm{def}})^2,
  \qquad0<\theta_h\leq1.
\end{equation}
If data oscillation dominates $\etaEq$, the same bulk criterion is applied
to $\operatorname{osc}_K^2$ instead.

\begin{algorithm}[Exact-certified adaptive LOD with cache batches]
\label{alg:exact-adaptive}
Given nested meshes $\mathcal T_H\preceq\mathcal T_h$, an oversampling
level $\ell$, parameters $\mu_0$, $\theta_H$, $\theta_h$,
$\tau_{\mathrm{loc}}$, $\tau_{\mathrm{tot}}$,
$q_{\mathrm{cert}}$, $m_{\mathrm{cert}}$, and a tolerance $\tolex$,
repeat the following steps.
\begin{enumerate}
\item \textbf{Admissibility and cached setup.}
Restore $\mathcal T_H\preceq\mathcal T_h$ and check
$\mu\leq\mu_0<1$.  If the resolution test fails, refine
$\mathcal T_H$ and apply the minimal fine-mesh nesting closure before any
corrector solve.  Assemble or update only the local corrector, Riesz, and
mixed factors affected by the last refinement.  Start a new cache batch if
$\mathcal T_h$ or $\ell$ changed.

\item \textbf{Solve and mark cheaply.}
Compute $\Qlr$, $\Qlrs$, solve \eqref{eq:PG}, evaluate
$\eta_{H,z}$ and $\eta_{H,T}$, and form a coarse D\"orfler set by
\eqref{eq:coarse-dorfler}.

\item \textbf{Decide whether to certify.}
If neither \eqref{eq:cert-trigger-reduction} nor
\eqref{eq:cert-trigger-delay} holds and termination is not being
considered, refine $\mathcal T_H$ on the coarse marked set, apply the
minimal nesting closure when necessary, and restart at Step~1.  Otherwise
compute $\etaEq$, $\underline\beta_\ell$, $\rhoP$, $\rhoA$,
$\betaInf$, and $\Uex$, and update the checkpoint records.

\item \textbf{Repair an uncertified corrector space.}
If $\epsA\geq1$, $\betaInf\leq0$, or
\eqref{eq:k-robust-acceptance} fails, compute $\deltaLocHat$.  If
$\deltaLocHat>\tau_{\mathrm{loc}}$, increase $\ell$ and restart at
Step~1.  Otherwise choose $\mathcal M_h$ by
\eqref{eq:fine-defect-dorfler}, refine $\mathcal T_h$ conformingly, and
restart at Step~1.

\item \textbf{Terminate or reduce the approximation error.}
If $\Uex\leq\tolex$, terminate with the certificate
\eqref{eq:exact-solution-upper}.  Otherwise refine $\mathcal T_H$ on the
coarse marked set.  If the prospective coarse mesh is not contained in
$\mathcal T_h$, refine the latter by the minimal nesting closure; this
structural update starts a new cache batch.  Continue with Step~1.
\end{enumerate}
\end{algorithm}

The algorithm gives a guaranteed stopping criterion whenever all
algebraic and geometric bounds are verified.  The D\"orfler rules for
$\eta_{H,T}$ and $\eta_{h,K}^{\mathrm{def}}$ are refinement strategies;
the present paper does not prove contraction or optimal complexity for
the coupled $H$--$h$--$\ell$ loop.  Such results would require estimator
reduction and quasi-orthogonality properties beyond the exact reliability
proved here.

A small fine-level reserve can be useful in practice.  If all meshes are
created by one bisection rule, one may monitor
\[
  g(\mathcal T_H,\mathcal T_h)
  :=\min_{T\in\mathcal T_H}
    \min_{\substack{K\in\mathcal T_h\\K\subseteq T}}
    (\operatorname{lev}K-\operatorname{lev}T)
\]
and apply a local fine-mesh closure before $g$ reaches zero.  This
structural safeguard reduces cache invalidations but has no role in the
error theorem; conforming bisection and its closure are discussed in
\cite{Stevenson2008}.

%----------------------------------------------------------------------
\subsection{Reuse and computational cost}
\label{sec:computational-realization}
%----------------------------------------------------------------------

Let
\[
  N_H:=\dim V_H,
  \qquad N_h:=\dim V_h.
\]
Denote by $n_{T,h}$ the dimension of a corrector problem on
$N^\ell(T)$, by $n_{z,h}$ the dimension of $W_{h,z}$, and by $n_{a,h}$
the number of mixed unknowns on a fine vertex patch.  Let
$\operatorname{Fact}(n)$ and $\operatorname{Solve}(n)$ denote one local
factorization and one back-solve.

On a corrector patch $D=N^\ell(T)$, the constraint $I_Hw_h=0$ can be
implemented by the saddle-point system
\begin{equation}
\label{eq:corrector-saddle-point}
  \begin{bmatrix}
    A_D&B_D^*\\ B_D&0
  \end{bmatrix}
  \begin{bmatrix}x\\\lambda\end{bmatrix}
  =\begin{bmatrix}r_T\\0\end{bmatrix},
\end{equation}
or by a null-space basis $Z_D$ and the reduced coercive matrix
$Z_D^*A_DZ_D$.  A complete setup costs schematically
\begin{equation}
\label{eq:corrector-cost}
  \sum_{T\in\mathcal T_H}\operatorname{Fact}(n_{T,h}),
\end{equation}
but within a cache batch only patches touched by a coarse refinement are
refactorized.  The patch problems are independent and parallelizable;
implementation strategies for LOD correctors and their reuse are discussed
in \cite{EngwerHenningMalqvistPeterseim2019}.

The residual indicator and localization diagnostic use the same positive
definite local $b_\kappa$ Gram matrices.  After factorization, one complete
$\eta_H$ evaluation has cost
\begin{equation}
\label{eq:etaH-cost}
  O(N_h)+\sum_{z\in\mathcal N_H}\operatorname{Solve}(n_{z,h}).
\end{equation}
This is the quantity evaluated at every coarse step.

The full-residual reconstruction and every action of $B_{\mathrm p}$ or
$B_{\mathrm a}$ use the same patchwise mixed matrix, with different
right-hand sides and the conjugated Robin sign in the adjoint case.  For
fixed polynomial degree and shape-regular meshes, one flux sweep costs
\begin{equation}
\label{eq:flux-sweep-cost}
  \sum_{a\in\mathcal N_h^{\mathrm{new}}}
      \operatorname{Fact}(n_{a,h})
  +\sum_{a\in\mathcal N_h}\operatorname{Solve}(n_{a,h})
  +O(N_h),
\end{equation}
which is linear in $N_h$ when the patch dimensions are uniformly bounded.
The first sum contains only newly created or structurally changed patches.

If the primal and adjoint generalized eigenproblems require
$m_{\mathrm p}$ and $m_{\mathrm a}$ matrix-free iterations, one exact
checkpoint costs approximately
\begin{equation}
\label{eq:checkpoint-cost}
  O\!\left((1+m_{\mathrm p}+m_{\mathrm a})N_h\right)
\end{equation}
local patch work, in addition to coarse linear algebra.  This can be
several tens of ordinary flux sweeps, so it should not be performed at
every coarse iteration.  Warm starts from the previous dominant
eigenvectors, block iterations, and parallel patch solves are important.
If $J$ exact checks are made during $M$ coarse iterations with $J\ll M$,
the average certification overhead is reduced by the factor $J/M$.
Moreover, $\rhoP$, $\rhoA$, and $\underline\beta_\ell$ depend on the
spaces and the wavenumber, but not on the forcing; they can be reused for
multiple right-hand sides as long as $\mathcal T_H$, $\mathcal T_h$, and
$\ell$ are unchanged.


%======================================================================
